\section{Bipartite Graph Check}
\label{sec:graph-bipartite}

\problemstatement{Given a graph, determine whether its nodes can be
colored using exactly \textbf{two} colors such that no two adjacent
nodes share the same color (such a graph is called \textbf{bipartite}).}

\begin{patternbox}
\emph{``Can this be split into two groups with no same-group edges''}
is a traversal that carries a \textbf{color} along: color the start
node $0$, every neighbor the \textbf{opposite} color, and so on
outward. The graph is bipartite exactly when this coloring never
forces a conflict -- an already-colored neighbor found to share the
current node's color.
\end{patternbox}

\subsection*{Why a conflict means \emph{not} bipartite}

If traversal reaches a node $v$ that already has a color, and that
color matches the current node $u$'s color, it means two different
paths from the start assigned $u$ and $v$ the \emph{same} color despite
$u$ and $v$ being directly connected -- which is only possible if the
cycle formed by those two paths plus the edge $u{-}v$ has \textbf{odd}
length (an even-length cycle alternates colors consistently all the
way around; an odd-length cycle cannot). So this check is, in effect,
testing for the absence of odd-length cycles.

\subsection*{Approach 1: BFS}

\begin{lstlisting}
#include <bits/stdc++.h>
using namespace std;

bool bfsCheckBipartite(int src, vector<vector<int>>& adj, vector<int>& color) {
    queue<int> q;
    color[src] = 0;
    q.push(src);

    while (!q.empty()) {
        int node = q.front();
        q.pop();

        for (int neighbor : adj[node]) {
            if (color[neighbor] == -1) {
                color[neighbor] = 1 - color[node];
                q.push(neighbor);
            } else if (color[neighbor] == color[node]) {
                return false;
            }
        }
    }
    return true;
}

bool isBipartite(int n, vector<vector<int>>& adj) {
    vector<int> color(n, -1);
    for (int i = 0; i < n; i++) {
        if (color[i] == -1 && !bfsCheckBipartite(i, adj, color)) {
            return false;
        }
    }
    return true;
}

int main() {
    int n = 4;
    vector<vector<int>> adj(n);
    auto addEdge = [&](int u, int v) { adj[u].push_back(v); adj[v].push_back(u); };
    addEdge(0, 1); addEdge(1, 2); addEdge(2, 3); addEdge(3, 0);

    cout << (isBipartite(n, adj) ? "true" : "false") << endl; // true (even cycle)
    return 0;
}
\end{lstlisting}

\begin{complexitybox}[Approach 1 Complexity]
\textbf{Time.} $O(V+E)$.
\textbf{Space.} $O(V)$ for \textit{color} and the queue.
\end{complexitybox}

\subsection*{Approach 2: DFS}

\begin{lstlisting}
#include <bits/stdc++.h>
using namespace std;

bool dfsCheckBipartite(int node, int c, vector<vector<int>>& adj, vector<int>& color) {
    color[node] = c;

    for (int neighbor : adj[node]) {
        if (color[neighbor] == -1) {
            if (!dfsCheckBipartite(neighbor, 1 - c, adj, color)) return false;
        } else if (color[neighbor] == c) {
            return false;
        }
    }
    return true;
}

bool isBipartiteDFS(int n, vector<vector<int>>& adj) {
    vector<int> color(n, -1);
    for (int i = 0; i < n; i++) {
        if (color[i] == -1 && !dfsCheckBipartite(i, 0, adj, color)) {
            return false;
        }
    }
    return true;
}

int main() {
    int n = 3;
    vector<vector<int>> adj(n);
    auto addEdge = [&](int u, int v) { adj[u].push_back(v); adj[v].push_back(u); };
    addEdge(0, 1); addEdge(1, 2); addEdge(2, 0);

    cout << (isBipartiteDFS(n, adj) ? "true" : "false") << endl; // false (odd cycle, triangle)
    return 0;
}
\end{lstlisting}

\subsection*{Dry run}

\dryrun{Take the triangle $0{-}1{-}2{-}0$ (a $3$-cycle), DFS from $0$.}

$\textit{dfs}(0,0)$: color $0$ with $0$. Neighbor $1$: uncolored,
recurse $\textit{dfs}(1,1)$: color $1$ with $1$. Neighbor $0$: colored
$0 \neq 1$, fine (it's the parent direction, naturally opposite).
Neighbor $2$: uncolored, recurse $\textit{dfs}(2,0)$: color $2$ with
$0$. Neighbor $1$: colored $1 \neq 0$, fine. Neighbor $0$: colored
$0$, and current color is also $0$ -- \textbf{conflict}, return
\texttt{false}. The odd cycle cannot be properly $2$-colored.

\begin{complexitybox}[Approach 2 Complexity]
\textbf{Time.} $O(V+E)$.
\textbf{Space.} $O(V)$ for \textit{color} and the recursion stack.
\end{complexitybox}

\begin{takeawaybox}
Carrying a small piece of alternating state (a color, a parity bit)
through a traversal is a general technique for detecting structural
properties tied to path \textbf{parity} -- odd cycles here, and the
same idea extends to any problem phrased as ``can these be split into
two groups with no same-group connections.''
\end{takeawaybox}
